% Generated draft: requires mathematical and source review.
\documentclass[11pt]{article}
\usepackage[margin=25mm]{geometry}
\usepackage{fontspec}
\setmainfont{Latin Modern Roman}
\usepackage{amsmath,amssymb,mathtools,unicode-math}
\setmathfont{Latin Modern Math}
\usepackage{enumitem,needspace,xurl,hyperref}
\setlength{\parindent}{0pt}
\setlength{\parskip}{5pt}
\setlength{\emergencystretch}{3em}
\newcommand{\N}{\mathbb{N}}
\newcommand{\Z}{\mathbb{Z}}
\newcommand{\Q}{\mathbb{Q}}
\newcommand{\R}{\mathbb{R}}
\newcommand{\C}{\mathbb{C}}
\newcommand{\F}{\mathbb{F}}
\newcommand{\A}{\mathbb{A}}
\newcommand{\PP}{\mathbb{P}}
\newcommand{\E}{\mathbb{E}}
\newcommand{\eps}{\varepsilon}
\newcommand{\dd}{\,\mathrm{d}}
\begin{document}
\section*{UnsolvedMath ID: 1100506}
\textbf{Source problem code:} catalogue entry
\subsection{Definitions and mathematical statement}
A group is virtually free if it has a free subgroup of finite index, allowing the trivial free group. For a discrete group $\Gamma\leq\operatorname{Isom}(\mathbb H^n)$ acting on real hyperbolic space of curvature $-1$, its critical exponent is
\[
 \delta(\Gamma)=\inf\left\{s\geq0:
 \sum_{\gamma\in\Gamma}e^{-s\,d(o,\gamma o)}<\infty\right\},
\]
with the usual infimum convention; it is independent of the basepoint $o\in\mathbb H^n$.

The transcribed question, attributed to Misha Kapovich, asks whether there is $\epsilon>0$ such that $\delta(G)<\epsilon$ implies that $G$ is virtually free. The displayed source question does not specify the class of groups, the meaning of $\delta$, finite generation, or dimensional dependence.

The supplied research assessment identifies the intended setting as finitely generated Kleinian groups and reports the following dimensional affirmative statement:
\[
 \forall n\geq2\ \exists D(n)\in(0,\tfrac12)\ \forall\Gamma\leq\operatorname{Isom}(\mathbb H^n),
\]
\[
 \bigl(\Gamma\text{ discrete and finitely generated},\
       \delta(\Gamma)<D(n)\bigr)
 \Longrightarrow\Gamma\text{ is virtually free}.
\]
This makes the assessment's dimension-dependent interpretation explicit; it does not assert a dimension-independent $\epsilon$ or an extension to groups outside that class. Whether the original list imposes further conventions must be checked in its surrounding context.
\subsection{Short English statement}
Does sufficiently small critical exponent force a finitely generated Kleinian group to be virtually free? The registry labels the question solved and attributes the affirmative dimension-dependent statement to Liu--Wang, Corollary 1.6.

The supplied, truncated summary describes the following mechanism. For $\delta<1$ it reports a Besson--Courtois--Gallot natural map with
$|\operatorname{Jac}_2F|\leq((1+\delta)/2)^2<1$, yielding a linear filling-area bound $4\ell(\mathcal C)/(1-\delta)$ for a null-homologous union of loops (Theorem 1.10). It reports cyclic parabolic subgroups, absence of parabolics if $\delta<1/2$, finitely many cusps and bounded geometry, and a characterization of convex cocompactness by properness of injectivity radius (Theorem 1.11). The required torsion, manifold, and other hypotheses for these subsidiary assertions are not supplied, so they are not asserted for every group in the displayed target.

The summary next attributes escaping closed geodesics to Kapovich--Liu and an infinite-descent/Margulis-tube argument to the Liu--Wang mechanism: two loxodromics of uniformly bounded displacement force a positive critical-exponent lower bound, contradicting $\delta<D(n,\kappa)$ and giving convex cocompactness (Theorem 1.2). The parameter $\kappa$ is not defined. It then invokes $\dim_H\Lambda(\Gamma)=\delta(\Gamma)<D(n)<1$, but ends mid-sentence before the concluding argument. These are attributed outline claims, not independently checked proofs or additional hypotheses silently inserted into the original question.
\subsection{Sources}
\begin{itemize}
\item UnsolvedMath permanent ID 1100506, source code catalogue entry, ``Questions in Geometric Group Theory --- Q 5.6.'' Registry curation and metadata: Ulam AI and UnsolvedMath Contributors, CC BY 4.0; updated 2026-08-06.
\item Bestvina, \emph{Questions in Geometric Group Theory} (source list dated 2004), Question 5.6, attributed to Misha Kapovich, PDF page 14: \url{https://www.math.utah.edu/~bestvina/eprints/questions-updated.pdf}. Supplied manual transcription accessed 2026-07-29. Original-list metadata presents an open item without independent current-status verification; status/release and rights review are NEEDS\_REVIEW.
\item The supplied resolving summary cites Liu--Wang, Corollary 1.6 and Theorems 1.2, 1.10, 1.11; it also names Besson--Courtois--Gallot and Kapovich--Liu. It supplies no complete titles, dates, or locators for these references; none are invented.
\end{itemize}
\end{document}
